\newpage
\section{Motivação}
\label{sec:motivacao}

O forró universitário difundiu-se nas universidades brasileiras ao longo das últimas três décadas, com presença documentada em campi de múltiplas regiões do país~\cite{havt2026dancar}. Apesar dessa popularização, o processo de ensino e aprendizagem dessas modalidades permanece fortemente ancorado na observação subjetiva de instrutores, que avaliam aspectos como postura, transferência de peso e adequação ao tempo musical com base exclusivamente em sua percepção visual e experiência pessoal \cite{porto2023ressignificando}. Essa dependência gera limitações importantes: a qualidade do feedback varia significativamente entre instrutores, o acesso a profissionais qualificados é restrito fora dos grandes centros urbanos, e o praticante raramente recebe métricas objetivas e replicáveis sobre sua evolução técnica ao longo do tempo. Ferramentas computacionais voltadas à análise de movimento têm avançado expressivamente em contextos esportivos e de reabilitação \cite{latyshev2024computer}, mas sistemas que integrem, de forma simultânea e acessível, a avaliação biomecânica do movimento e a validação rítmica por metrônomo ainda são escassos, especialmente voltados a modalidades de dança de salão \cite{THARATIPYAKUL2024e36589}.

A literatura registra evidências de benefícios da dança de salão à saúde física e mental, incluindo melhora da coordenação motora, do equilíbrio, da capacidade cardiorrespiratória e da autoimagem corporal dos praticantes \cite{fonseca2012, rodrigues2011, paiva2019danca}. Nesse sentido, o desenvolvimento de uma ferramenta capaz de oferecer feedback técnico individualizado, sem a necessidade de equipamentos especializados ou supervisão constante de um instrutor, possui aplicação prática direta: democratizar o acesso a recursos de aprendizagem que hoje são privilégio de quem tem condições financeiras de frequentar aulas particulares ou estúdios com infraestrutura avançada de captura de movimento.

Do ponto de vista científico, este trabalho se conecta a uma tendência maior de democratização de tecnologias de análise de desempenho. Recursos antes restritos a laboratórios de biomecânica e centros de treinamento de alto rendimento -- como captura de movimento, estimativa de pose em tempo real e inferência de modelos de aprendizado de máquina -- tornaram-se viáveis em dispositivos móveis de uso cotidiano graças a avanços recentes em arquiteturas de inferência otimizadas para hardware embarcado \cite{lee2019ondevice, li2020edgeai}. O presente trabalho explora essa convergência ao propor um sistema que transfere para dispositivos móveis de consumo capacidades antes restritas a laboratórios de análise de movimento.

Esse movimento é comprovado por dados recentes de mercado, que evidenciam a velocidade com que soluções de inteligência artificial vêm sendo incorporadas ao universo do esporte, da atividade física e do bem-estar. Estima-se que o mercado global de inteligência artificial aplicada ao esporte, avaliado em aproximadamente US\$ 10,82 bilhões em 2025, deverá atingir cerca de US\$ 60,78 bilhões até 2034, com uma taxa composta de crescimento anual (CAGR) superior a 21\% \cite{precedence2025aisports}. No contexto brasileiro, o mercado de fitness digital -- que engloba aplicativos de treino, acompanhamento de desempenho e tecnologias vestíveis -- já é estimado em torno de US\$ 1,32 bilhão em 2025, com projeções globais que apontam o setor de fitness digital se aproximando de US\$ 98,7 bilhões, impulsionado por startups que utilizam IA para personalização de treinos \cite{saudedigital2025fitness}. Esses números indicam não apenas uma oportunidade de mercado, mas também uma mudança de comportamento dos usuários, que estão cada vez mais dispostos a utilizar seus próprios smartphones como ferramentas de acompanhamento de atividades físicas -- cenário favorável à adoção de um aplicativo de análise de dança baseado em visão computacional embarcada.

Por fim, este tema possui também uma motivação pessoal e acadêmica para os autores. Como praticante de forró universitário e estudante de Engenharia de Computação na UNIFEI, os autores identificaram, em suas próprias trajetórias como alunos de dança, a dificuldade de obter feedback técnico consistente sobre aspectos como postura e adequação ao ritmo musical, especialmente em etapas iniciais de aprendizagem. Essa vivência pessoal, somada ao interesse acadêmico em inteligência artificial aplicada a problemas do mundo real, motivou o desenvolvimento do Dança~AI como objeto de estudo na interseção entre visão computacional e ensino de dança.